\documentclass[sigconf]{acmart}

% CSC 510 course report rather than an ACM publication.
\setcopyright{none}
\settopmatter{printacmref=false}
\renewcommand\footnotetextcopyrightpermission[1]{}
\pagestyle{plain}
\raggedbottom

\usepackage{array}
\usepackage{booktabs}
\usepackage{graphicx}
\usepackage{hyperref}
\usepackage{placeins}
\newcolumntype{L}[1]{>{\raggedright\arraybackslash}p{#1}}
\setlength{\emergencystretch}{1.5em}

\newcommand{\pending}[1]{\textit{[#1]}}

\title{Project 1b}
\acmConference[CSC 510]{CSC 510: Software Engineering}{September 2026}{Raleigh, NC}
\acmYear{2026}

\author{Abigail Close}
\affiliation{%
  \institution{North Carolina State University}
  \city{Raleigh}
  \state{North Carolina}
  \country{USA}}
\email{aeclose2@ncsu.edu}

\author{Aditya Mahajan}
\affiliation{%
  \institution{North Carolina State University}
  \city{Raleigh}
  \state{North Carolina}
  \country{USA}}
\email{amahaja8@ncsu.edu}

\author{Satwi Shah}
\affiliation{%
  \institution{North Carolina State University}
  \city{Raleigh}
  \state{North Carolina}
  \country{USA}}
\email{stshah2@ncsu.edu}

\author{Supreme Constantine}
\affiliation{%
  \institution{North Carolina State University}
  \city{Raleigh}
  \state{North Carolina}
  \country{USA}}
\email{spconsta@ncsu.edu}

\begin{document}

\begin{abstract}
WolfBite is an inherited Flutter/Firebase shopping assistant for WIC
participants and caregivers. Project 1a established the system's existing
behavior and limitations through traceable use cases and independent testing.
In Project 1b, we combined analyses from multiple LLMs with verified competitor
sources, participant evidence, Project 1a engineering results, feasibility
review, and relevant regulations, standards, and licenses to assess the market,
identify product opportunities, and consider whether to continue with WolfBite
or pivot. The evidence indicates a medium-confidence opportunity to help
shoppers respond when an item that appears eligible is rejected at checkout.
The team chose to continue developing WolfBite. Project 2 will pursue the
finalized one-month milestones: clearer nutrition explanations, bounded
basket-wide suggestions, a more usable shopping interface, and rule-based help
for simulated checkout-rejection scenarios, including a baseline comparison of
the checkout-help workflow. This plan extends the inherited product while
keeping uncertain eligibility and checkout information explicit.
\end{abstract}

\maketitle

\section{Introduction}

Project 1a established the ``Before'' baseline for the inherited WolfBite
application through code-traceable use cases and automated testing. Project 1b
used that baseline to determine whether the team should extend WolfBite or
pivot. To support this decision, the team compared four LLM analyses with
verified competitor sources, participant evidence, feasibility findings, and
relevant technical and regulatory material. This work produced a market survey,
a documented product gap, a mission and stakeholder analysis, scoped
engineering milestones, and an assessment of both continuation and pivot
options.

The team decided to continue with WolfBite and selected the finalized ``Now''
milestones as the Project 2 plan. These milestones cover evaluation of the
checkout-help concept, clearer nutrition explanations, basket-wide suggestions,
interface improvements, and bounded help for simulated item rejections. The
following sections present the evidence and constraints behind that decision
and define how the proposed version will improve on the Project 1a system.

\section{D1: Market Survey}

\subsection{Method}

We ran the shared market-survey Prompt~1 with four LLM analysts: Codex,
Terra, Gemini, and local Ollama, to generate rival candidates. We retained a
rival only when at least two models named it or one model supplied a live,
relevant URL. We then checked official product sources and independently
audited the three strongest rivals. Each final rival independently satisfies
the live-URL qualification rule through an official product source.

\subsection{Verified Rivals}

Table~\ref{tab:rivals} summarizes the three verified rivals.

\begin{table*}[t]
  \caption{D1 validated WIC rivals.}
  \label{tab:rivals}
  \centering
  \small
  \begin{tabular}{L{2.3cm}L{3.4cm}L{4.3cm}L{4.8cm}}
    \toprule
    Rival & Evidence it is real & What it does & Relevant limitation \\
    \midrule
    WICShopper &
    Official App Store listing and vendor purchase-history documentation
    \cite{wicshopper-appstore,wicshopper-redemption-history} &
    Shows benefits, scans eligibility, provides food guidance, and records
    item/category purchase history in supported agencies &
    No documented workflow connects a scan/checkout mismatch to its likely
    cause and recovery steps. \\
    \addlinespace
    myWIC Mosaic &
    Official App Store listing \cite{mywic-mosaic-appstore} &
    Shows benefits, tracks purchases, scans products, and supports shopping,
    appointments, and certification &
    No documented item-level explanation and recovery workflow for an
    unexpectedly uncovered item. \\
    \addlinespace
    Bnft &
    Official App Store listing \cite{bnft-appstore} &
    Shows balances, scans benefit eligibility, manages cards, and provides
    transaction history &
    No documented item-level reconciliation or explanation for conflicting
    scan and checkout results. \\
    \bottomrule
  \end{tabular}
\end{table*}

\subsection{Selected Gap and Audience}

\textbf{Gap:} Our feature audit found no documented integrated, item-level
explanation-and-recovery workflow in WICShopper, myWIC Mosaic, or Bnft for
cases where an app indicates that a product is eligible but checkout rejects
it. Such a workflow could connect the rejected item, the participant's current
benefit or food category, and purchase evidence to likely causes and
actionable next steps.

\textbf{Audience:} WIC caregivers and participants shopping for their
households, especially those managing young children while an item they
expected to be covered is rejected at checkout.

This need is supported by participant research reporting difficulty
identifying WIC-eligible foods and stigma during shopping and checkout
\cite{fiedler2026,chauvenet2019}. WIC caregivers also rank balance checking and
barcode scanning among their highest-priority app features \cite{weber2021}.
Together, this evidence suggests that eligibility uncertainty during shopping
remains consequential for the target audience.

We also evaluated two alternative gaps. We rejected nutrition comparison
because participant evidence for that need was comparatively weak. We rejected
broader receipt reconciliation because all three rivals already provide some
form of purchase or transaction history. The narrower explanation-and-recovery
gap therefore has stronger support from both the competitor audit and
participant evidence.

\subsection{Counterevidence and Scope}

The conclusion has \textbf{medium confidence}. WICShopper already provides
item-level purchase history in some agencies, myWIC tracks purchases, and Bnft
provides transaction history. Public documentation may also omit
program-specific features. In one month, the team can build and test a
WolfBite prototype using fixed mismatch scenarios, rule-based likely causes,
and recovery guidance. Without state benefit-system data, the prototype must
not claim to know an authoritative point-of-sale rejection reason.

\section{D2: Mission, Impact, and Stakeholders}

\subsection{Mission Statement}

WIC participants and caregivers can face uncertainty when a product appears
eligible in an app but is rejected at checkout, a setting in which participant
research has documented difficulty identifying eligible foods and stigma
during shopping and checkout \cite{fiedler2026,chauvenet2019}. That conflict
can add confusion and friction during a household shopping trip. We propose
extending WolfBite with an item-level explanation-and-recovery workflow that
connects the rejected item's UPC and package details, the participant's current
benefit or food category, and available eligibility or purchase evidence. For
a bounded set of mocked mismatch scenarios, the workflow will distinguish app
evidence from facts that only a retailer or WIC agency can confirm, then show
rule-based likely causes and actionable next steps. This guidance is intended
to help participants and caregivers choose a next action after an unexpected
rejection while keeping uncertain information visible.

\subsection{Measurable Claim and M0 Evaluation}

\textbf{M0 measurable claim:} During Project 2 M0, we will test whether the
explanation-and-recovery workflow produces a task-success rate of at least
80\% and at least 20 percentage points higher than the baseline workflow. A
trial succeeds when a test user selects the prepared scenario's reference next
step within 60 seconds without facilitator help. The baseline is a control
version of the current WolfBite flow that shows the pre-checkout eligibility
result and a mocked checkout rejection but provides no explanation or recovery
guidance. The new and baseline flows will use the same counterbalanced,
synthetic mismatch scenarios, so the evaluation requires neither live WIC
data nor retailer point-of-sale access. These values are project targets to be
measured at M0, not achieved results.

\subsection{Stakeholders}

The four Prompt~5 outputs used different labels for related roles. To avoid
inflating the list with synonyms, Table~\ref{tab:stakeholders} includes a
stakeholder group when there is a direct equivalent between LLMs.

\begin{table*}[t]
  \caption{Overlapping Prompt 5 stakeholders and requirements consequences.}
  \label{tab:stakeholders}
  \centering
  \small
  \setlength{\defaultaddspace}{0.3em}
  \begin{tabular}{L{3.0cm}L{5.0cm}L{7.0cm}}
    \toprule
    Stakeholder & Need or fear & Design response \\
    \midrule
    WIC participants, caregivers, and household representatives &
    Misleading eligibility guidance, uncertainty after a rejection, or the
    wrong household state &
    Identify the active household context and show evidence limits, likely
    causes, and a short recovery action for the selected item. \\
    \addlinespace
    Infants, children, and pregnant or postpartum benefit recipients &
    Eligibility or nutrition guidance could be mistaken for advice about what
    is healthy or appropriate &
    Keep coverage guidance separate from health claims and represent missing
    information as unknown rather than zero or safe. \\
    \addlinespace
    WIC clinic staff, program administrators, and agencies &
    Extra support burden or a prototype mistaken for official WIC authority &
    Show data source and freshness, and route unresolved cases to the
    appropriate agency contact. \\
    \addlinespace
    Cashiers, authorized retailers, and store managers &
    Checkout delay, disputes, or incorrect UPC and package information when
    the app and register disagree &
    Treat the register outcome as authoritative, provide concise recovery
    guidance, and make no payment or point-of-sale integration claim. \\
    \addlinespace
    State WIC IT teams, EBT processors, and card-system vendors &
    Unsupported integrations, security risks, or simulated balances presented
    as official data &
    Use mock adapters and simulated data, document the integration boundary,
    and make no live balance or payment claim. \\
    \addlinespace
    Approved-product-list stewards, manufacturers, and product-data providers &
    Stale, copied, misattributed, or misidentified UPC, package, and category
    records &
    Retain source/version provenance and test stale, duplicate, and conflicting
    records before presenting a likely cause. \\
    \addlinespace
    Privacy, security, and legal reviewers &
    Exposure of household, benefit, receipt, location, or diagnostic data and
    liability from unsupported claims &
    Minimize stored data, exclude card credentials, use privacy-safe logs, and
    state the prototype's limits at decision points. \\
    \addlinespace
    Program regulators and compliance personnel &
    Misrepresented eligibility, unnecessary data collection, or exclusion of
    protected users &
    Trace eligibility statements to their source and test explicit privacy,
    accessibility, and nondiscrimination requirements. \\
    \addlinespace
    Participants with disabilities or limited English proficiency, and their
    advocates &
    Camera-only controls, color-only status, small targets, or unclear language
    could block the recovery task &
    Preserve manual entry, pair color and icons with text, use semantic labels
    and plain language, and test text scaling and assistive navigation. \\
    \addlinespace
    Receipt OCR and other outside data-service providers &
    Sensitive uploads, unsupported formats, or use outside service terms &
    Keep OCR outside the core recovery path; if retained, require consent,
    validate inputs, minimize retention, and provide a mock service. \\
    \addlinespace
    Application developers and maintainers &
    Unreproducible failures, sensitive logs, or one-month scope growth &
    Use a bounded rule set, deterministic fixtures, and privacy-safe
    diagnostics. \\
    \addlinespace
    Project sponsors, funders, and procurement decision-makers &
    Misleading public-benefit claims, hidden integration costs, or a prototype
    that cannot be maintained &
    Define a narrow pilot, measurable outcomes, and explicit external
    dependencies rather than promising statewide deployment. \\
    \bottomrule
  \end{tabular}
\end{table*}
\FloatBarrier

\section{D3: Engineering Plan}

\subsection{Before: Project 1a}

\begin{itemize}
  \item \pending{Add three to five completed, evidenced milestones.}
\end{itemize}

\subsection{Now: One-Month Build and Test}

Our one-month goal is to make WolfBite a more complete, usable shopping app.
We will improve nutrition explanations, add an optional basket-wide suggestion
mode, and make the main screens easier to use. A small checkout-help feature
retains D2's explanation-and-recovery goal; the nutrition improvements support
the shopping experience rather than establish a new verified market gap.
Benefit information and checkout rejection scenarios will remain simulated.

\begin{itemize}
  \item \textbf{M0: Check whether checkout help is useful.}
  Compare the current flow with the checkout-help feature using the same
  prepared rejection scenarios. Keep D2's target: at least 80\% of trials
  choose the reference next step within 60 seconds without help, with success
  at least 20 percentage points above the baseline. Build a script that scores
  trial logs and runs in CI; check it with known passing and failing examples.
  Synthetic logs test the script, not real usability. Report pilot results
  separately if volunteers are available, and record a re-plan decision if
  measured results miss the target.

  \item \textbf{M1: Better nutrition information and food-choice explainers.}
  Extend the existing nutrition panel and badges with units, serving or
  reference amounts, and short explanations of the values behind a label or
  comparison. Show available sodium, sugar, fiber, protein, and fat information,
  including tradeoffs between options rather than a single unexplained
  ``healthier'' score. Reuse the nutrition utility and product records, preserve
  missing values as unknown, and retain the measurement basis in saved basket
  items. Check educational wording against cited nutrition sources during
  implementation. Unit and widget tests must verify displayed values and units,
  explanation rules, comparable serving bases, and missing-data cases; unknown
  values must never earn a favorable badge or comparison.

  \item \textbf{M2: Basket-wide suggestions for more balanced food choices.}
  Add an optional ``Basket balance'' mode alongside individual substitutions.
  Use basket quantities and M1's nutrition data to summarize food-group variety
  and rank possible swaps by their effect on the whole basket. Start with a
  small mock catalog with known package amounts and food-group mappings, and
  explicit demonstration targets for sodium, added sugar, and fiber. Prioritize
  only comparisons with known, compatible units and package amounts; flag
  incomplete totals and skip suggestions that depend on missing measurements.
  Rank candidates against
  the basket's largest target gap, rather than ranking each food in isolation;
  show before/after values and any tradeoffs for up to three suggestions.
  Reuse the APL category lookup, restrict swaps to the same benefit category
  and compatible package size, and recheck simulated allowances before applying
  one shopper-confirmed swap at a time. Recalculate after each change. Tests
  must show that the same item can receive different suggestions in different
  baskets, quantities affect results, and previews leave the basket unchanged.
  Also test confirmed swaps, balance consistency, missing measurement data,
  and no suitable candidate. This is a shopping aid with illustrative targets,
  not a personalized diet plan or a claim that a basket represents daily intake.

  \item \textbf{M3: A clearer, more complete shopping interface.}
  Improve the existing scan, basket, and benefits screens with consistent
  typography, spacing, product cards, and clearly labeled primary actions.
  Make nutrition explanations and the two suggestion modes easy to find;
  distinguish covered and shopper-paid items with text as well as color.
  Add useful loading, empty, error, and retry states. Widget tests and a manual
  walkthrough must cover scan-to-basket navigation, opening an explainer,
  previewing a basket suggestion, and checking benefits. Verify phone and
  desktop layouts, enlarged text, and screen-reader labels without clipped
  content or inaccessible controls. Existing quantity and navigation tests
  must continue to pass.

  \item \textbf{M4: Add straightforward help for a rejected item.}
  Add an action on a basket item that opens a short explanation and next step
  for prepared package-size, category-balance, and stale-information scenarios.
  Use simple rules over mock item and benefit records, with an ``unable to
  determine'' response when evidence is missing. Test each explanation and
  the full item-to-help flow, including returning to the basket without
  changing its contents. Label the advice as possible causes, not an official
  checkout decision; use this feature in M0's comparison.
\end{itemize}

Build the M0 script and M1's measurement handling in week 1, implement and test
basket suggestions in weeks 2--3, and develop the interface and bounded
checkout-help cases alongside them. Reserve week 4 for integration and the
evaluation. Keep basket suggestions to one-swap previews over the mock catalog;
multi-item optimization and personalized meal planning are outside this month.
Add tests alongside each change, run relevant Project 1a regressions in CI,
and document how another team can run the mock demo and extend the catalog,
explanation rules, basket targets, and checkout-help cases.

\subsection{Future: After the One-Month Project}

\begin{itemize}
  \item \pending{Add three to five credible future milestones.}
\end{itemize}

\subsection{Support Material}

Table~\ref{tab:support} separates material that governs the WIC program from
technical and accessibility guidance that constrains this prototype and from
licenses that govern redistribution.  The design consequences are requirements
for the Project~2 explanation-and-recovery workflow, not claims that WolfBite
is a WIC agency, retailer system, or EBT processor.

\begin{table*}[t]
  \caption{Support material and the Project~2 design checks it creates.}
  \label{tab:support}
  \centering
  \footnotesize
  \begin{tabular}{L{3.8cm}L{5.1cm}L{6.1cm}}
    \toprule
    Support material and source & What it says that matters here & Effect on WolfBite's design \\
    \midrule
    \textbf{Federal regulation: 7 CFR Part 246, especially
    \S~246.10(b) and (e)(12).} USDA Food and Nutrition Service (FNS);
    \href{https://www.ecfr.gov/current/title-7/subtitle-B/chapter-II/subchapter-A/part-246/section-246.10}{eCFR text} and
    \href{https://www.fns.usda.gov/wic/food-packages/regulatory-requirements}{FNS regulatory summary}. &
    The federal rule supplies minimum food and package criteria.  It also
    permits a state agency to add criteria and requires the agency to identify
    acceptable brands and package sizes.  FNS therefore cautions that a state
    need not authorize every food that meets the federal requirements. &
    An eligibility result must identify its jurisdiction, exact UPC/package,
    source, and source date.  Meeting only federal criteria will be labeled as
    general guidance, never as a guarantee of purchase.  The workflow will
    distinguish that evidence from state authorization and from the retailer's
    actual checkout result. \\
    \addlinespace[0.2em]
    \textbf{State domain guidance: North Carolina WIC Authorized Product List
    (APL).} NC Department of Health and Human Services (NCDHHS);
    \href{https://www.ncdhhs.gov/divisions/child-and-family-well-being/community-nutrition-services-section/wic/vendors/nc-wic-authorized-product-list-apl}{APL and UPC-submission page} and
    \href{https://www.ncdhhs.gov/wic-vendor-manual-0}{WIC Vendor Manual}. &
    NCDHHS maintains an APL by UPC; products can be added, changed, or removed,
    and vendor systems must receive APL updates.  A product that appears to fit
    a federal category is not thereby present in North Carolina's current APL. &
    The bounded NC scenarios will use a dated APL fixture keyed by UPC and
    package details.  A missing, changed, or stale APL record may be shown only
    as a likely cause.  Recovery will link to the current NC list and agency
    help or suggest checking an authorized alternative, rather than declaring
    the item approved or diagnosing the register. \\
    \addlinespace[0.2em]
    \textbf{Technical/domain guidance: WIC EBT Technical Implementation Guide
    and Operating Rules.} USDA FNS;
    \href{https://www.fns.usda.gov/wic/ebt/technical-implementation-guide-operating-rules}{official resource page}. &
    These resources standardize WIC EBT online purchase messages and file
    handling for state agencies, authorized vendors, processors, and other EBT
    stakeholders.  They describe the transaction ecosystem; they do not turn
    a shopping assistant into an EBT system. &
    Fixtures and explanations will keep agency, APL, benefit, and retailer
    evidence separate.  Project~2 will use mock adapters and no card number,
    PIN, live purchase message, or balance mutation.  It will say ``likely
    cause'' when only the state or retailer transaction system could supply an
    authoritative result. \\
    \addlinespace[0.2em]
    \textbf{Accessibility standard/design target: WCAG 2.2 Level AA.}
    World Wide Web Consortium (W3C);
    \href{https://www.w3.org/TR/WCAG22/}{WCAG 2.2} and
    \href{https://www.w3.org/WAI/standards-guidelines/mobile/}{W3C mobile guidance}. &
    Relevant success criteria prohibit color as the only means of conveying
    information (1.4.1), set minimum contrast (1.4.3) and target-size/spacing
    rules (2.5.8), require identified errors and known corrections
    (3.3.1/3.3.3), and keep repeated help consistent (3.2.6).  W3C explains how
    WCAG guidance can also inform native mobile applications. &
    Approval, rejection, and uncertainty will use text plus an icon, not color
    alone.  The flow will use semantic labels, readable contrast, scalable
    text, adequately sized or spaced controls, a manual-UPC alternative to the
    camera, a specific error and next action, and consistently placed agency
    help.  Tests will cover text scaling and assistive navigation. \\
    \addlinespace[0.2em]
    \textbf{Application license: MIT License.} The inherited WolfBite
    \texttt{Project3/LICENSE.md}; the same license is visible in the
    \href{https://github.com/SuyeshJadhav/CSC510_G19/blob/main/Project2/LICENSE.md}{upstream repository}. &
    MIT permits use, modification, and redistribution, provided the copyright
    and permission notice remain in copies or substantial portions; it also
    disclaims warranties. &
    We will keep the inherited copyright and MIT text with redistributed
    WolfBite code and will not describe the prototype as warranted or officially
    approved. \\
    \addlinespace[0.2em]
    \textbf{Selected dependency licenses.} The declared packages are in
    \texttt{Project3/pubspec.yaml}; official package records confirm
    \href{https://pub.dev/packages/provider/license}{provider (MIT)},
    \href{https://pub.dev/packages/mobile_scanner/license}{mobile\_scanner (BSD 3-Clause)}, and
    \href{https://pub.dev/packages/firebase_core/license}{firebase\_core (BSD 3-Clause)}. &
    MIT requires its copyright and permission notice.  BSD 3-Clause requires
    its copyright, conditions, and disclaimer in source or binary distributions
    as applicable and prohibits using contributor names for endorsement. &
    Releases will preserve applicable third-party notices.  Before adding an
    APL, accessibility, or scanning package, the team will record the resolved
    version and license, check compatibility, and reject any dependency whose
    redistribution duties the project cannot meet. \\
    \bottomrule
  \end{tabular}
\end{table*}

\section{D5: LLM Analysis and Verification}

\subsection{Models, Prompts, and Evidence}

The four model/runner pairings were GPT-5.6 Sol through Codex desktop,
collected by Abigail Close; GPT-5.6 Terra through ChatGPT, collected by Supreme
Constantine; Gemini 3.1 Pro (High) through Antigravity, collected by Aditya
Mahajan; and Llama 3 8B through local Ollama, run by Satwi Shah. All four
completed Prompt~1, the required market survey. The team's exact P01 keeper is
preserved in \path{Project1b_Work/prompts/keeper_prompts.md}; the Codex and
Gemini records embed that version, while the Terra and local records preserve
the similar earlier wording they actually received. Complete raw transcripts
are stored by model under \path{Project1b_Work/model_outputs/}. Rival
normalization and verification are separate from those raw outputs in
\path{Project1b_Work/rivals/validated_rivals.md}.

The local P01 run produced ten candidates without live browsing. Verification
retained its normalized WICShopper entry as a direct rival, classified several
real products as merely adjacent, and rejected or left unverified the remaining
unsupported candidates. The run is therefore useful evidence about local-model
limitations, but not independent confirmation of its links or feature claims.

The finalized P10 keeper and Codex run are preserved in the same prompt and
model-output directories. \pending{TODO --- awaiting team input: run the
finalized P10 with Terra and Gemini and rerun it with local Ollama; the existing
local P10 used the superseded milestone set.}

\subsection{Prompt-by-Model Comparison}

\begin{table*}[t]
  \caption{Comparison of the market-survey and red-team prompts.}
  \label{tab:model-comparison}
  \centering
  \small
  \begin{tabular}{L{2.2cm}L{2.0cm}L{3.6cm}L{3.6cm}L{3.6cm}}
    \toprule
    Model & Prompt & Most useful result & Weakness or disagreement & Team adjudication \\
    \midrule
    Codex & P01 &
    Returned ten candidates with live evidence URLs and warned that an omitted
    feature is not proof of a market gap. &
    Mixed platform-level products with state-specific applications; detailed
    feature and price claims still required checking. &
    Verify every product and narrow the strongest direct set to WICShopper,
    myWIC Mosaic, and Bnft. \\
    \addlinespace
    Terra & P01 &
    Focused on direct state WIC applications and identified geographic
    fragmentation. &
    Supplied source labels rather than the requested URL strings. &
    Retain candidates corroborated elsewhere; exclude the three Terra-only
    candidates that lacked a qualifying URL. \\
    \addlinespace
    Gemini & P01 &
    Corroborated the three strongest rivals and several state applications with
    relevant links. &
    Returned seven rather than ten rows and included weakness claims not
    established by its linked product pages. &
    Verify product existence independently and discard the unsupported
    weakness claims. \\
    \addlinespace
    Local Ollama & P01 &
    Identified WICShopper and exposed how a non-browsing model expands toward
    adjacent grocery products. &
    Included dead or irrelevant links, non-WIC products, and unsupported
    receipt-scanning claims. &
    Normalize WIC Shopper to WICShopper; do not count the remaining entries as
    strong direct rivals. \\
    \addlinespace
    \multicolumn{5}{l}{\pending{TODO --- awaiting team input: add comparable P10 rows for all four models.}} \\
    \bottomrule
  \end{tabular}
\end{table*}

\subsection{Wrong Outputs and How We Caught Them}

The team checked candidate identity, URL relevance, product scope, and feature
claims separately instead of accepting model prose as evidence. This process
caught the following concrete problems in P01:

\begin{itemize}
  \item Gemini returned seven products, not the requested ten. We counted the
  raw rows and evaluated those seven rather than inventing replacements. Its
  claims about scanner accuracy, frequent login errors, and reinstalling were
  not established by the supplied product listings, so those claims were not
  used as market findings.
  \item Terra placed source descriptions such as ``Washington State DOH'' in
  the evidence field rather than URLs. Under the qualification rule, its
  uncorroborated Indiana WIC, Wisconsin MyWIC, and SC WIC Mobile App entries
  were excluded.
  \item The local model supplied dead, irrelevant, or unsupported links for
  FreshConnect, a generic ``WIC Mobile App,'' StoreWise, and Grocery IQ. Live
  URL and scope checks removed them from the qualifying direct-rival set.
  Ibotta, Fetch Rewards, Checkout 51, Out of Milk, and AnyList were verified as
  adjacent rebate or grocery tools, not direct WIC competitors.
  \item The local output attributed receipt scanning to Out of Milk and
  AnyList, but the reviewed product pages did not support those claims. The
  team discarded them rather than using them to justify a receipt feature.
  \item An initial team note also called the local run's WICShopper URL
  apparently fabricated and claimed zero overlap with Terra. Name normalization
  and a live redirect check showed that ``WIC Shopper'' was WICShopper. The
  final audit therefore credited all four models and retained it as the
  strongest direct rival while preserving the incorrect note for transparency.
\end{itemize}

These corrections prevented unsupported products and weaknesses from becoming
requirements. They also narrowed the final D1 claim from broad nutrition or
receipt novelty to a medium-confidence explanation-and-recovery gap.

\subsection{Most and Least Useful Prompts}

P01 was useful as a discovery prompt: its four runs supplied rival candidates,
revealed the state-by-state fragmentation of WIC applications, and provided
names for independent verification. It was not sufficient as market proof.
Even the strongest run required feature checking, and the local run showed how
confident but unverified links and adjacent products could distort the rival
set. The team therefore used P01 to generate candidates, not to establish the
final gap. \pending{TODO --- awaiting team input: compare the completed
common-version P10 outputs before naming the final most and least useful
prompts.}

\subsection{Model Strengths and Weaknesses}

Based on P01, the models contributed different strengths and failure modes:

\begin{itemize}
  \item \textbf{Codex} supplied the most complete set of live links and an
  explicit warning against inferring feature absence, but its broad list mixed
  market platforms with state-specific implementations.
  \item \textbf{Terra} stayed close to direct WIC workflows and exposed program
  fragmentation, but its source labels did not satisfy the requested URL
  evidence format.
  \item \textbf{Gemini} provided relevant corroborating products and links,
  but stopped at seven entries and attached unverified failure-frequency claims
  to several rivals.
  \item \textbf{Local Ollama} usefully demonstrated the risk of running market
  discovery without browsing, but most of its list was unsupported or adjacent
  to WIC rather than directly competitive.
\end{itemize}

This disagreement improved the audit: the browsing models supplied stronger
direct candidates, while the local model made the need for URL and scope
verification unmistakable. \pending{TODO --- awaiting team input: revise these
lines if the common-version P10 runs reveal additional model-specific strengths
or weaknesses.}

\section{Conclusion}

\pending{Restate the verified gap, product decision, measurable benefit, and
why the one-month plan is both ambitious and testable.}

\bibliographystyle{ACM-Reference-Format}
\bibliography{references}

\end{document}
